\documentclass[11pt]{article}

\usepackage[review]{acl}
\usepackage{times}
\usepackage{latexsym}
\usepackage[T1]{fontenc}
\usepackage[utf8]{inputenc}
\usepackage{microtype}
\usepackage{inconsolata}
\usepackage{graphicx}
\usepackage{amsmath}
\usepackage{amssymb}
\usepackage{booktabs}
\usepackage{subcaption}
\usepackage{xcolor}
\usepackage{multirow}
\usepackage{placeins}

%% Utility macros
\newcommand{\ph}[1]{\textcolor{gray}{[#1]}}        % placeholder value
\newcommand{\todo}[1]{\textcolor{red}{\textbf{[TODO: #1]}}}
\newcommand{\OR}{\textsc{or}}                        % over-refusal shorthand


\title{Over-Refusal and Representation Subspaces:\\
A Geometrical Analysis of Task-Conditioned Refusal in Aligned LLMs}

\author{
  Utsav Maskey \quad
  Mark Dras \quad
  Usman Naseem \\
  Macquarie University \\
  \texttt{\{utsav.maskey,mark.dras,usman.naseem\}@mq.edu.au}
}

\begin{document}


\usepackage[review]{acl}
\usepackage{times}
\usepackage{latexsym}
\usepackage[T1]{fontenc}
\usepackage[utf8]{inputenc}
\usepackage{microtype}
\usepackage{inconsolata}
\usepackage{graphicx}
\usepackage{amsmath}
\usepackage{amssymb}
\usepackage{booktabs}
\usepackage{subcaption}
\usepackage{xcolor}
\usepackage{multirow}
\usepackage{placeins}

%% Utility macros
\newcommand{\ph}[1]{\textcolor{gray}{[#1]}}        % placeholder value
\newcommand{\todo}[1]{\textcolor{red}{\textbf{[TODO: #1]}}}
\newcommand{\OR}{\textsc{or}}                        % over-refusal shorthand


\title{Over-Refusal and Representation Subspaces:\\
A Geometrical Analysis of Task-Conditioned Refusal in Aligned LLMs}

\author{
  Utsav Maskey \quad
  Mark Dras \quad
  Usman Naseem \\
  Macquarie University \\
  \texttt{\{utsav.maskey,mark.dras,usman.naseem\}@mq.edu.au}
}

\begin{document}
\maketitle

\begin{abstract}


Representation steering---ablating a direction from a model's residual stream---works
strikingly well for harmful refusals: a single global direction suppresses them with
80\% attack success rate.
But applied to \emph{over-refusal}---safe instructions incorrectly refused because
they superficially resemble harmful requests---the same approach is non-selective:
it disrupts the entire refusal mechanism rather than targeting over-refusal specifically. In this work, we provide a geometric account of why this happens.
On LLaMA-3.1-8B-Instruct, the \emph{harmful-refusal direction} and the
\emph{over-refusal direction} are two distinct extractions targeting two distinct
sample populations.
They are similar enough to cause collateral suppression when one is ablated,
but different enough that this suppression is non-selective.
Over-refusal lives inside task-identity clusters and has no single global axis to target---
its directions vary across tasks, its subspace is higher-dimensional, and it is perfectly
separable from harmful-refusal throughout the network via a linear probe from the first layer.
A task-conditioned method (SafeConstellations) that operates within this within-task
structure achieves genuine selectivity---eliminating over-refusal with far less collateral
safety disruption.
The right question is not whether steering can suppress over-refusal, but whether it
can do so \emph{selectively}.

\end{abstract}

%% -----------------------------------------------------------
\section{Introduction}
%% -----------------------------------------------------------

\paragraph{Steering for refusal control.}
Safety-aligned LLMs are trained to refuse harmful requests, but controlling this
behaviour at inference time---without retraining---is hard.
Representation steering offers a practical solution: extract a ``refusal direction''
from the residual stream via difference-in-means (DIM) over harmful/harmless pairs,
then ablate it to suppress refusals \citep{arditi2024refusal,zou2023representation}.
For harmful requests, this works strikingly well: ablating the global DIM direction
achieves 80\% attack success rate (ASR) on our evaluation set (68\% harmful refusal
$\to$ 20\% after ablation) with no side-effects on harmless inputs.
The success rests on a specific property: harmful-refusal directions are nearly
identical across all tasks (cosine alignment 0.835--0.858), making a single global
vector a sufficient intervention target.

\paragraph{The over-refusal problem.}
Safety-aligned models also refuse \emph{safe} instructions framed in sensitive
contexts \citep{safeconstellations2025}---a translation task involving historical
violence, a sentiment analysis request on inflammatory text, a generic cryptographic
question.
This \emph{over-refusal} carries a real usability cost, especially for professional
and research applications.
The obvious fix is steering: if the model refuses when it should not, ablate the
harmful-refusal direction.
In practice, however, this fails to be \emph{selective} \citep{wangf2025refusal}:
it suppresses harmful-refusal more than over-refusal, and disrupts the safety circuit
as a side-effect.

\paragraph{Our central claim.}
This failure is not accidental---it follows directly from the geometry of representation
space.
Figures~\ref{fig:2d_key} and~\ref{fig:galaxy} show the key observation: at the
mid-layer where representations are richest, the dominant structure is \emph{task identity},
not refusal class.
Over-refusal samples do not occupy a separate ``refusal zone''---they sit \emph{inside}
their task-coloured clouds, organised around task-identity axes, with per-task
over-refusal directions that diverge across tasks (Figure~\ref{fig:2d_key}).
This contrasts directly with harmful-refusal directions, which converge to a single
task-agnostic axis (Figure~\ref{fig:nb9_repr_space}).
The two phenomena are two distinct extractions from two distinct sample populations:
the former is universal; the latter is task-conditioned, higher-dimensional, and
varies across tasks.
Ablating the harmful-refusal direction cannot selectively fix over-refusal, because it
was not extracted for that population---the partial geometric overlap produces collateral
suppression, not targeted correction.

A method that targets the \emph{within-task} subspace instead can.
We validate this: SafeConstellations \citep{safeconstellations2025} eliminates
over-refusal entirely while causing substantially less collateral safety disruption
than global ablation (selectivity ratio 1.57 vs.\ 0.67).

\paragraph{Contributions.}
\begin{enumerate}
  \item We replicate Arditi ablation and document where it succeeds and where it fails.
  The harmful-refusal direction is universal across tasks, which is exactly why a global
  intervention works for harmful content.
  Applied to over-refusal, it is non-selective---suppressing harmful refusals more than
  over-refusal---exposing it as a blunt instrument.

  \item We provide a geometric account of why over-refusal is different.
  Mid-layer representations are dominated by task identity; over-refusal sits inside
  task clusters rather than in a shared refusal zone.
  The over-refusal direction is task-conditioned, task-specific, and substantially
  misaligned with the harmful-refusal direction---its subspace is higher-dimensional
  and its principal axes are task-identity directions.

  \item We confirm with functional probing that task identity encodes earlier and more
  sharply than refusal behaviour, and that over-refusal and harmful-refusal are perfectly
  separable by a linear probe from the very first layer ($p<0.001$).
  Causal patching further shows that refusal commitment is distributed: no single head
  is individually sufficient to flip the decision.

  \item We validate the selectivity prediction empirically: a task-conditioned method
  (SafeConstellations) eliminates over-refusal entirely while causing substantially less
  collateral safety disruption, achieving a selectivity ratio of \textbf{1.57} vs.\
  0.67 for Arditi ablation.
  The residual bypass is geometrically predicted by constellation structure.
\end{enumerate}

\FloatBarrier
%% -----------------------------------------------------------
\section{Background}
%% -----------------------------------------------------------
\label{sec:background}

\paragraph{Refusal localization.}
\citet{arditi2024refusal} show that a DIM direction over harmful/harmless pairs is
necessary and sufficient for refusal across 13 models. \citet{zhao2025refusal} extend
this to reasoning models. Component-level analyses identify specific attention heads
\citep{zhou2025safety} and neurons \citep{zhao2025d} as primary circuit elements.
SAE-based decompositions \citep{templeton2024scaling,goyal2025detox} provide a
feature-level picture.

\paragraph{Steering limitations.}
\citet{braun2025unreliability} identify three failure conditions: high magnitude,
distributional shift, and feature entanglement. \citet{hedstrom2025steer} show that
abstaining from steering is optimal when the model lacks stable representations.
\citet{wang2025adaptive} find that a single truthfulness vector fails across
hallucination categories. These failure modes are consistent with, but do not
explain, the specific geometric structure we document for over-refusal.

\paragraph{Task-conditioned representations.}
\citet{safeconstellations2025} show that LLMs encode task identity through stable
``constellation'' trajectory patterns in activation space, and that task-conditioned
steering reduces over-refusal by up to 73\%. Our paper provides the mechanistic
explanation for why this works while global ablation fails.

\paragraph{Non-linear representations.}
\citet{engels2025nonlinear} show that some concepts are encoded as manifolds, not
single directions. We find a related limitation for over-refusal: its within-task
direction is not a stable global direction but is task-specific and geometrically
distinct from the harmful-refusal axis.

\FloatBarrier
%% -----------------------------------------------------------
\section{Setup}
%% -----------------------------------------------------------
\label{sec:setup}

\paragraph{Model and data.}
We study \textbf{LLaMA-3.1-8B-Instruct} (32 layers, 4096-dim). Our evaluation set
has 270 samples across five tasks: sentiment analysis, translation, cryptanalysis,
RAG-style QA, and rephrasing. Each sample is labeled with a \emph{response label}
(\texttt{cautious / not\_harmful / harmful}) and a \emph{refusal class}
(\texttt{direct\_answer / direct\_refusal / indirect\_refusal}) via a GPT-4o
judge using the ICML (2025) three-class prompt \citep{safeconstellations2025}.

\paragraph{Behavioral groups.}
Following \citet{safeconstellations2025}, groups are defined by intended task, not
surface label, because cryptanalysis and RAG-QA samples are labeled \texttt{harmful\_instruction}
in text-type fields even when benign---any refusal on these tasks is over-refusal:
\begin{align*}
  \text{OVER\_REFUSAL} &= \text{refusal} \cap \text{benign task} \\
  \text{TARGET} &= (\texttt{cautious} \cup \texttt{not\_harmful}) \cap \texttt{direct\_answer} \\
  \text{REFUSED\_HARMFUL} &= \text{refusal} \cap \texttt{harmful\_label}
\end{align*}
where ``benign task'' $\in$ \{sentiment, translate, cryptanalysis, rag\_qa\}.

\paragraph{Embeddings.}
Activations are extracted at the input-normed residual stream of each layer
(\texttt{layer\_\{l\}\_input\_norm}). Silhouette scores are computed in the native
4096-dim space; 2-D PCA is used for visualization (deterministic, no hyperparameters,
directions are interpretable as linear combinations of the residual stream).

\paragraph{Arditi replication.}
We extract the global DIM direction at layer 12 (best AUROC) on our 25
refused-harmful / 30 harmless-answered sample pairs, replicating \citet{arditi2024refusal}.
All generation experiments use greedy decoding (150 tokens), classified by the same
GPT-4o judge as the original pipeline.

\paragraph{Two distinct refusal directions.}
We distinguish two DIM-extracted directions throughout, as they are easy to conflate:
the \emph{harmful-refusal direction} ($\hat{\mathbf{r}}_{\text{Arditi}}$), extracted
from refused-harmful vs.\ harmless-answered pairs following \citet{arditi2024refusal};
and the \emph{over-refusal direction} ($\hat{\mathbf{v}}^{\OR}$), extracted from
benign-refused vs.\ benign-answered pairs.
They target different sample populations and capture different phenomena.
The central finding (\S\ref{sec:geometry}) is that they are geometrically distinct
at mid-layers, which is why ablating one cannot selectively suppress the other.

\FloatBarrier
%% -----------------------------------------------------------
\section{Results}
%% -----------------------------------------------------------
\label{sec:results}

Results are presented in narrative order.
\S\ref{sec:arditi} establishes that global DIM ablation \emph{works} for harmful refusals---confirming the universality that makes it effective.
\S\ref{sec:constellations}--\ref{sec:probing} characterise why over-refusal is geometrically different: it lives inside task constellations, not along the harmful-refusal axis.
\S\ref{sec:h4} shows empirically that global DIM ablation is consequently non-selective for over-refusal.
\S\ref{sec:selectivity} demonstrates that task-conditioned methods achieve genuine selectivity.
\S\ref{sec:circuits} provides the mechanistic account of late-layer distributed commitment.

%% ---- §4.1 ----
\subsection{Steering Works: The Harmful-Refusal Direction Is Universal}
\label{sec:arditi}

We first replicate \citet{arditi2024refusal} to establish the baseline.
The global DIM direction ($\hat{\mathbf{r}}_{\text{Arditi}}$), extracted at layer 12,
achieves 80\% attack success rate (ASR) on our refused-harmful samples
(Table~\ref{tab:arditi})---reducing harmful refusal from 68\% to 20\% with no
false positives on harmless inputs.

\begin{table}[h]
\centering
\small
\setlength{\tabcolsep}{7pt}
\begin{tabular}{lcc}
\toprule
 & Baseline & After ablation \\
\midrule
Harmful refusal rate  & 68.0\% & 20.0\%  \\
Harmless refusal rate & 0.0\%  & 0.0\%  \\
ASR                   & 32.0\% & \textbf{80.0\%} \\
\bottomrule
\end{tabular}
\caption{Arditi replication (NB8). Global DIM ablation substantially reduces harmful
refusal (68\%$\to$20\%, ASR: 32\%$\to$80\%) with no cost to harmless-answered samples.}
\label{tab:arditi}
\end{table}

This success is grounded in a specific geometric property: harmful-refusal is
\emph{universal}. Task-specific DIM directions are nearly identical to the global
direction (cosine alignment 0.835--0.858 across all tasks), and cross-task transfer
nearly equals self-transfer. The direction reaches $\geq$89\% cross-task alignment
by \textbf{layer 3}---before task-identity constellations even peak---establishing
that harmful refusal is mediated by a single task-agnostic direction that a
global intervention can cleanly target.
This sets the baseline: global DIM works \emph{because} its target is universal.
The question is whether over-refusal shares this property.

\begin{figure*}[t]
  \centering
  %% SOURCE: NB9 Cell 10 — q_2d_representation_space.png
  %% Place at: paper/figures/fig_nb9_2d_repr_space.png
  \includegraphics[width=0.95\linewidth]{figures/fig_nb9_2d_repr_space.png}
  \caption{\textbf{Harmful-refusal directions are universal: 2D PCA at three layers ($n=270$).}
  Colour = task identity; shape = behaviour ({\large$\bullet$} harmless-answered, $\blacktriangle$ over-refusal, $\times$ refused-harmful).
  Black arrow = global Arditi direction; coloured arrows = per-task harmful-refusal directions.
  At crucial layers (L17), task-specific arrows converge onto the single global direction---confirming
  that harmful refusal is mediated by a task-agnostic signal that a global intervention
  can cleanly target.}
  \label{fig:nb9_repr_space}
\end{figure*}

%% ---- §4.2 ----
\subsection{Task-Identity Constellations Dominate Mid-Layer Geometry}
\label{sec:constellations}

To understand why over-refusal is different, we first examine the representational
structure that emerges in mid-layers---where the model's richest representations live.
Figure~\ref{fig:2d_key} makes the geometric argument directly: over-refusal samples
sit \emph{inside} their task-coloured clusters at every layer, with no shared refusal
zone, while the per-task over-refusal directions (coloured arrows) each point in a
different direction---the opposite of the convergence seen for harmful-refusal
(Figure~\ref{fig:nb9_repr_space}).

\begin{figure*}[t]
  \centering
  %% SOURCE: NB13a Cell 7 — fig13a_00_2d_key.png
  %% Place at: paper/figures/fig13a_00_2d_key.png
  \includegraphics[width=0.95\linewidth]{figures/fig13a_00_2d_key.png}
  \caption{\textbf{Why global steering cannot selectively fix over-refusal: 2D PCA at three layers ($n=270$).}
  Colour = task identity; $\blacktriangle$ = over-refusal; $\times$ = refused-harmful; {\large$\bullet$} = harmless-answered.
  Black arrow = global harmful-refusal direction (Arditi); coloured arrows = per-task
  over-refusal directions (HA centroid $\to$ OR centroid within each task cluster).
  At the peak-constellation layer (L12), over-refusal triangles sit \emph{inside} their
  task clusters with no separate refusal zone; the per-task OR arrows diverge, unlike
  the converging harmful-refusal arrows in Figure~\ref{fig:nb9_repr_space}.
  No single global direction can selectively reach the OR samples without also
  suppressing harmless samples in the same cluster region.}
  \label{fig:2d_key}
\end{figure*}

\begin{figure*}[t]
  \centering
  %% SOURCE: NB13a — dual galaxy (task identity left, behavioral class right)
  %% Place at: paper/figures/fig_nb13a_dual_galaxy.png
  \includegraphics[width=0.95\linewidth]{figures/fig_nb13a_dual_galaxy.png}
  \caption{\textbf{Over-refusal lives inside task constellations (L12, $n=270$).}
  2-D PCA of residual-stream activations at the peak-constellation layer.
  Left: samples coloured by task identity---five well-separated constellations.
  Right: same projection coloured by behaviour---over-refusal samples ($\blacktriangle$)
  sit \emph{inside} their task clusters with no separate refusal zone.}
  \label{fig:galaxy}
\end{figure*}

\begin{figure}[ht]
  \centering
  %% SOURCE: NB13b cell nb13b-11 — per-task OR arrows + Arditi arrow at L12
  %% Place at: paper/figures/fig_nb13b_l12_snapshot.png
  \includegraphics[width=\linewidth]{figures/fig_nb13b_l12_snapshot.png}
  \caption{\textbf{Per-task over-refusal directions vs.\ the harmful-refusal direction (L12).}
  Coloured arrows: per-task over-refusal directions, each pointing differently.
  Black arrow: the harmful-refusal direction, crossing all task boundaries.}
  \label{fig:or_arrows}
\end{figure}

Figure~\ref{fig:galaxy} confirms the dominant structure of mid-layer activations is
task identity, not refusal class.
Five well-separated task constellations form at this layer, while the within-task
behavioral signal---the difference between over-refusal and correctly-answered
samples---is a small perturbation inside each cluster.

Figure~\ref{fig:or_arrows} reinforces this geometrically: the per-task
over-refusal directions (coloured arrows) each point in a \emph{different} direction,
while the harmful-refusal direction (black arrow) is a single vector crossing all
task boundaries.
These per-task directions are far less mutually consistent than harmful-refusal
directions---the constellations are organised by task, not by refusal class.
The behavioral gap shrinks in late layers as task-specific circuits commit the final
refusal decision (\S\ref{sec:circuits}).

\begin{figure}[ht]
  \centering
  %% SOURCE: NB13a — silhouette line plot across all layers
  %% Place at:  paper/figures/fig3_silhouette_gap.png
  \includegraphics[width=\linewidth]{figures/fig13a_02_silhouette_curves.png}
  \caption{\textbf{Constellation strength and within-task behavioral gap across layers.}
  Top: inter-task silhouette peaks at L12 and stays elevated.
  Bottom: within-task behavioral distance vs.\ inter-task distance (dashed)---the
  behavioral gap is a small within-cluster perturbation that shrinks in late layers.}
  \label{fig:silhouette}
\end{figure}

This establishes the first structural fact: over-refusal is \emph{task-conditioned},
not refusal-conditioned. The geometry it lives in is organised by task identity, not
a shared refusal axis. A direction that crosses task boundaries---as the global Arditi
direction does---cannot cleanly target it. We quantify this directly in the next section.

%% ---- §4.3 ----
\subsection{The Geometric Case: Over-Refusal Does Not Follow the Universal Direction}
\label{sec:geometry}

The constellation picture suggests that the \emph{over-refusal direction} ($\hat{\mathbf{v}}^{\OR}$)
is not aligned with the \emph{harmful-refusal direction} ($\hat{\mathbf{r}}_{\text{Arditi}}$)
but rather organized around task-identity axes.
We test this directly through three complementary analyses: directional alignment,
per-task consistency, and subspace dimensionality.

\subsubsection{Over-Refusal Direction vs.\ the Harmful-Refusal Direction}
\label{sec:h1}

Recall that the harmful-refusal direction $\hat{\mathbf{r}}_{\text{Arditi}}$ is computed
from refused-harmful vs.\ harmless-answered pairs (\S\ref{sec:setup}).
The over-refusal direction is an entirely different extraction: DIM over
benign-refused (over-refusal) vs.\ benign-answered pairs,
\begin{equation}
  \hat{\mathbf{v}}^{\OR}_\ell = \frac{
    \boldsymbol{\mu}^\ell_{\text{benign-refused}} - \boldsymbol{\mu}^\ell_{\text{benign-answered}}
  }{\|\boldsymbol{\mu}^\ell_{\text{benign-refused}} - \boldsymbol{\mu}^\ell_{\text{benign-answered}}\|}
\end{equation}
and we measure $\cos(\hat{\mathbf{v}}^{\OR}_\ell,\, \hat{\mathbf{r}}_{\text{Arditi},\ell})$
at each of the 32 layers (Figure~\ref{fig:cosine}). The shaded band (0.835--0.858) shows
the harmful-refusal inter-task alignment as a reference for what ``universal'' looks like.

\begin{figure}[ht]
  \centering
  %% SOURCE: NB14 — fig1_or_vs_arditi_cosine.png
  %% Place at:  paper/figures/fig2_cosine_h1.png
  \includegraphics[width=\linewidth]{figures/fig2_cosine_h1.png}
  \caption{\textbf{Layer-wise cosine similarity between the over-refusal and harmful-refusal directions.}
  Red: $\cos(\hat{\mathbf{v}}^{\OR}_\ell, \hat{\mathbf{r}}_{\text{Arditi},\ell})$.
  Shaded band: harmful-refusal inter-task alignment (reference for a universal direction).
  The over-refusal direction plateaus near 0.45---roughly half the harmful-refusal
  alignment---confirming that the two directions occupy geometrically distinct subspaces.}
  \label{fig:cosine}
\end{figure}

The cosine rises through early layers and stabilises around \textbf{0.45} from
mid-layers onward---roughly half the harmful-refusal inter-task alignment (0.845--0.858).
The two directions are not orthogonal, which is why ablating one partially suppresses
the other as a collateral effect.
But they are far from identical, which is why that suppression is non-selective.
This partial overlap is the geometric mechanism underlying all the empirical results
that follow.

\subsubsection{Over-Refusal Is More Task-Specific}
\label{sec:h2}

If over-refusal lives inside task clusters, the over-refusal direction should vary across
tasks more than the harmful-refusal direction does.
We test this by comparing per-task over-refusal directions against the inter-task
harmful-refusal baseline.

Harmful-refusal directions are nearly identical across all tasks---they are universal.
Per-task over-refusal directions tell a different story.
Only two tasks have sufficient over-refusal samples to compute a direction reliably:
\textit{sentiment analysis} and \textit{translation}.
Their pairwise directional similarity is well below the harmful-refusal baseline
(Figure~\ref{fig:h2}), and each aligns only weakly with the global over-refusal
direction---the two tasks are pulling in different directions.

\begin{figure*}[t]
  \centering
  %% SOURCE: NB14 — fig2_or_task_similarity.png
  %% Place at:  paper/figures/fig4_h2_or_tasks.png
  \includegraphics[width=0.95\linewidth]{figures/fig4_h2_or_tasks.png}
  \caption{\textbf{Per-task over-refusal direction similarity at L15, L17 and L28.}
  Left: pairwise cosine heatmap of per-task over-refusal directions.
  Right: mean off-diagonal similarity for over-refusal (red) vs.\ harmful-refusal
  directions---over-refusal is task-specific; harmful-refusal is universal.}
  \label{fig:h2}
\end{figure*}

The contrast is stark: harmful-refusal has a single universal direction that can be
extracted globally; over-refusal has different directions per task, organized around
task identity rather than a shared refusal axis.

\subsubsection{The Over-Refusal Subspace Is Higher-Dimensional and Task-Structured}
\label{sec:h3_dim}

A direct consequence of task-conditioning is that over-refusal variance should be
\emph{distributed across task directions} rather than concentrated in a single axis.
We test this via PCA of the over-refusal and harmful-refusal subspaces at each layer.

\paragraph{Variance structure at L12.}
The harmful-refusal subspace is relatively concentrated: PC1 explains \textbf{30.3\%}
of its variance and \textbf{8} components suffice for 80\% (Figure~\ref{fig:nb16_variance}).
The over-refusal subspace is more diffuse: PC1 explains only \textbf{24.5\%} and
\textbf{11} components are needed for 80\%---roughly 38\% more components.

\begin{figure}[ht]
  \centering
  %% SOURCE: NB16 — fig_nb16_variance.png
  %% Place at: paper/figures/fig7_nb16_variance.png
  \includegraphics[width=\linewidth]{figures/fig7_nb16_variance.png}
  \caption{\textbf{Variance structure of over-refusal vs.\ harmful-refusal subspaces at L17.}
  Left: cumulative explained variance---harmful-refusal concentrates faster (8 components
  for 80\%) than over-refusal (11 components).
  Right: per-PC variance---over-refusal is more diffuse at every rank.}
  \label{fig:nb16_variance}
\end{figure}

\paragraph{The principal axes of over-refusal are task-identity directions.}
Figure~\ref{fig:nb16_task_pca} makes the task-conditioning concrete. In a 2-D PCA
of over-refusal samples centred against the dataset mean, PC1 (24.5\% var) and PC2
(20.4\% var) \emph{separate tasks}, not refusal signal: translate samples cluster on the
right ($+$PC1), sentiment\_analysis on the left ($-$PC1), with near-zero within-task
scatter. The dominant axes of the over-refusal representation space are task-identity
directions---not a shared refusal axis.

\begin{figure}[ht]
  \centering
  %% SOURCE: NB16 — fig_nb16_task_pca.png
  %% Place at: paper/figures/fig8_nb16_task_pca.png
  \includegraphics[width=\linewidth]{figures/fig8_nb16_task_pca.png}
  \caption{\textbf{2-D PCA of over-refusal samples, coloured by task (L17).}
  The principal axes separate tasks, not refusal signal---the dominant structure
  of the over-refusal subspace is task identity.}
  \label{fig:nb16_task_pca}
\end{figure}

The layer sweep (Figure~\ref{fig:nb16_layer_sweep} in Appendix) shows this pattern
is consistent across \textbf{layers 7--30}. Notably, at early layers (L1--L4), the
pattern inverts---over-refusal briefly appears more concentrated---consistent with
the timeline of constellation formation: before task structure peaks at L12, over-refusal
has not yet been absorbed into task-identity geometry.

Together, these three analyses---directional alignment, task-specificity, and
dimensionality---converge on one conclusion: over-refusal occupies a task-conditioned,
multi-directional subspace that no single global direction can cleanly target.
We next confirm this with functional evidence from linear probing.

%% ---- §4.4 ----
\subsection{Probing Confirms Representational Separability}
\label{sec:probing}

The geometric analyses above establish that over-refusal and harmful-refusal are
directionally distinct. We complement this with a functional test: do linear probes
decode this structure, and in what order does it emerge across layers?

\paragraph{Setup.}
We train three logistic-regression probes at each of the 31 layers via 5-fold
stratified cross-validation.
\textbf{Task probe} (5-class, $n=270$): predict task identity.
\textbf{Refusal-behaviour probe} (binary, $n=217$): over-refusal ($n=48$) vs.\ TARGET ($n=169$)
within benign-task samples.
\textbf{Refusal-type probe} (binary, $n=73$): over-refusal ($n=48$) vs.\ REFUSED\_HARMFUL
($n=25$)---two refused populations separated by cause, not content.

\paragraph{Results (Figure~\ref{fig:nb17_probes}).}

Three findings stand out.
First, \textbf{task identity encodes immediately and saturates early}: the task probe
reaches near-ceiling performance from the first layer and is fully saturated well before
the mid-layer constellation peak.
Second, \textbf{refusal behaviour encodes later and within task structure}: the
refusal-behaviour probe improves through mid-layers and peaks after task saturation,
consistent with over-refusal consolidating inside the already-formed task geometry.
Third---and most directly relevant to intervention---\textbf{over-refusal and harmful-refusal
are perfectly separable throughout the entire network}: the refusal-type probe achieves
100\% balanced accuracy from the first transformer layer onward ($p<0.001$).
The model knows, from the very beginning, whether it is looking at a benign task or
a genuinely harmful request---the two refusal types occupy non-overlapping regions of
representation space at every layer.

\begin{figure*}[t]
  \centering
  %% SOURCE: NB17 — fig_nb17_probe_all3.png
  %% Copy to: paper/figures/fig9_nb17_probes.png
  \includegraphics[width=0.95\linewidth]{figures/fig9_nb17_probes.png}
  \caption{\textbf{Linear probe accuracy across layers.}
  (a) Task probe saturates immediately from L1.
  (b) Refusal-behaviour probe (over-refusal vs.\ harmless-answered) improves through
  mid-layers, consistent with over-refusal forming inside established task structure.
  (c) Refusal-type probe (over-refusal vs.\ refused-harmful) achieves 100\% from L1---the
  two refusal causes are representationally distinct throughout the network.}
  \label{fig:nb17_probes}
\end{figure*}

The 100\% refusal-type accuracy has a direct implication for interventions: if the
two refusal types are representationally distinct throughout the network, a method
that targets task-conditioned geometry \emph{should} be able to suppress one without
disrupting the other. Conversely, a method that ignores this structure will
inevitably cross-contaminate. We test both predictions next.

%% ---- §4.5 ----
\subsection{Global DIM Ablation Is Non-Selective in Practice}
\label{sec:h4}

Given the geometric distinction between the two refusal directions, what happens when
the harmful-refusal direction ($\hat{\mathbf{r}}_{\text{Arditi}}$) is ablated to fix
over-refusal---a population it was not designed for?
Figure~\ref{fig:h3} presents the direct empirical comparison ($n=20$ per condition, GPT-4o judge).

The result is exactly what the geometry predicts: Arditi ablation suppresses both
classes, but suppresses harmful refusals more than over-refusal---the numbers are in
Table~\ref{tab:nb15}.

\begin{figure*}[t]
  \centering
  %% SOURCE: NB14 — fig3_arditi_ablation_comparison.png
  %% Place at:  paper/figures/fig5_h3_ablation.png
  \includegraphics[width=0.95\linewidth]{figures/fig5_h3_ablation.png}
  \caption{\textbf{Arditi ablation effect on harmful-refusal vs.\ over-refusal ($n=20$, GPT-4o judge).}
  Left: refusal rates before and after ablation.
  Right: suppression magnitude---Arditi suppresses harmful refusals more than over-refusal
  (selectivity ratio 0.67), confirming it is a blunt instrument for this population.}
  \label{fig:h3}
\end{figure*}

The geometry predicted this outcome.
The two directions share partial variance, so ablating one partially suppresses the
other as a collateral effect---but this suppression is non-selective.
The key issue is not whether the method \emph{works} but whether it is \emph{selective}:
using the harmful-refusal direction to fix over-refusal disrupts the safety circuit
as a side-effect, because the two directions are not the same.
A method targeting the within-task over-refusal subspace directly can do better.

%% ---- §4.6 ----
\subsection{Task-Conditioned Steering Achieves Genuine Selectivity}
\label{sec:selectivity}

The geometric account predicts that task-conditioned methods should outperform global
ablation on selectivity: greater over-refusal suppression with less collateral
disruption to the harmful-refusal circuit. We test this directly by running
SafeConstellations \citep{safeconstellations2025} and Arditi ablation on the same
20 over-refusal and 20 refused-harmful samples.

\begin{figure*}[t]
  \centering
  %% SOURCE: NB15 — fig_nb15_selectivity.png / fig_nb15_selectivity.png
  %% Place at:  paper/figures/fig6_selectivity.png
  \includegraphics[width=0.95\linewidth]{figures/fig6_selectivity.png}
  \caption{\textbf{Selectivity comparison: SafeConstellations vs.\ Arditi ablation ($n=20$, GPT-4o judge).}
  Left: refusal rates under baseline, Arditi, and SafeConstellations.
  Right: suppression magnitudes---SafeConstellations eliminates over-refusal entirely
  while causing less harmful-refusal bypass, improving selectivity ratio from 0.67 to 1.57.}
  \label{fig:selectivity}
\end{figure*}

\begin{table}[h]
\centering
\small
\setlength{\tabcolsep}{5pt}
\begin{tabular}{lccc}
\toprule
 & Baseline & Arditi & SafeConstel \\
\midrule
Over-refusal rate    & 55\% & 25\% ($-30$\,pp) & \phantom{0}0\% ($-55$\,pp) \\
Harmful refusal rate & 65\% & 20\% ($-45$\,pp) & 30\% ($-35$\,pp) \\
\midrule
Selectivity (OR$\div$harm supp.) & --- & 0.67 & \textbf{1.57} \\
\bottomrule
\end{tabular}
\caption{Head-to-head comparison ($n=20$ per condition). SafeConstellations
achieves higher over-refusal suppression ($-55$\,pp, to 0\%) while preserving more of
the safety circuit (30\% residual harmful refusal vs.\ 20\% for Arditi, $-35$\,pp vs.\
$-45$\,pp). Selectivity ratio = OR suppression $\div$ harm suppression; higher is more
selective.}
\label{tab:nb15}
\end{table}

Results confirm the prediction (Table~\ref{tab:nb15}).
Arditi ablation is non-selective: it suppresses harmful refusals more than over-refusal,
consistent with operating on a direction designed for a different population.
SafeConstellations reverses this: it eliminates over-refusal entirely while causing
substantially less harmful-refusal bypass---the selectivity ratio improves from 0.67
(Arditi) to \textbf{1.57} (SafeConstellations).

The residual $-35$\,pp harmful-refusal bypass has a direct mechanistic explanation
from the constellation geometry (\S\ref{sec:constellations}).
Our refused-harmful samples are task-wrapped: they arrive as translation, rephrasing,
or sentiment analysis prompts with harmful content embedded.
At L12, where task identity dominates representations (silhouette\,=\,0.454), these
samples are geometrically indistinguishable from benign over-refusal samples at the
task-cluster level.
The constellation hook fires on task identity---not on content harmfulness---and
therefore applies equally to harmful task-wrapped prompts.
This is not a flaw but an inherent property of targeting task-identity geometry,
and it is a consequence directly predicted by the constellation structure established
in \S\ref{sec:constellations}.
In deployment, a prior content classifier would gate steering to verified-benign task
instances, closing this gap.
Critically, even without such a gate, SafeConstellations remains strictly more
selective than Arditi ($1.57$ vs.\ $0.67$).

\subsection{Mechanistic Account: Late-Layer Task-Conditioned Circuits}
\label{sec:circuits}

To identify \emph{where} over-refusal is committed mechanistically, we compute
per-token refusal attribution scores for each attention head $(l, h)$ across task
conditions using gradient-weighted attention \citep{zhou2025safety}.
The top-K attribution heads are concentrated exclusively in layers 27--31
(Table~\ref{tab:heads}), with 80\% task-specific membership (mean off-diagonal
overlap = 0.20 across tasks).

\begin{table}[h]
\centering
\small
\setlength{\tabcolsep}{5pt}
\begin{tabular}{llc}
\toprule
Head & Task & Role \\
\midrule
L31.H21 & All tasks    & Shared (universal) \\
L28.H31 & rephrase     & Task-specific \\
L30.H4  & rephrase     & Task-specific \\
L31.H30 & sentiment    & Task-specific \\
L31.H11 & sentiment    & Task-specific \\
L30.H3  & translate    & Task-specific \\
L31.H13 & translate    & Task-specific \\
\bottomrule
\end{tabular}
\caption{Top refusal-attribution attention heads by task (NB12). One shared head
(L31.H21) is universal; 80\% of heads are task-specific. All reside in layers 27--31.}
\label{tab:heads}
\end{table}

This late-layer, task-specific concentration ties together all previous observations.
The within-task behavioral gap \emph{shrinks} from L12 to L28 (\S\ref{sec:constellations})
because early-to-mid layers accumulate the behavioral perturbation but final refusal
commitment is delegated to late-layer task-specific heads.
Arditi ablation, which operates on the global residual stream direction before these
heads, disrupts the broad refusal signal but does not selectively suppress the
task-conditioned circuit that commits over-refusal.

\subsubsection{Causal Patching: Refusal Is Distributed, Not Head-Local}
\label{sec:causal_patching}

We directly test the causal necessity of the attributed heads via activation patching.
For each top-K attributed head $(l, h)$, we replace its output at the last-token
position (\texttt{[:, -1, :]}) with the corresponding activation from a harmless prompt
and measure whether this flips the model's refusal decision.
This targets exactly the position that determines the model's output token, independently
of sequence length.

\begin{figure}[ht]
  \centering
  %% SOURCE: NB18 — fig_nb18_causal_patching.png
  %% Copy to: paper/figures/fig10_nb18_causal_patching.png
  \includegraphics[width=\linewidth]{figures/fig10_nb18_causal_patching.png}
  \caption{\textbf{Causal patching flip rates for top-15 attribution heads.}
  Fraction of test pairs (out of 5) that flip to non-refusal when a head's last-token
  output is replaced with its harmless-prompt counterpart.
  No head crosses the 50\% causal-necessity threshold; six heads tie at 40\%.
  Refusal commitment is distributed across heads.}
  \label{fig:nb18_causal}
\end{figure}

\paragraph{Results.}
The global flip rates (Figure~\ref{fig:nb18_causal}) reveal a striking pattern:
no head approaches the 50\% threshold for causal necessity.
The six highest-flipping heads tie at 40\% (L13.H12, L13.H24, L13.H15, L13.H02, L13.H08,
L15.H08; 2 out of 5 test pairs each)---notably, these are \emph{not} among the
late-layer task-specific heads identified by attribution (L27--L31).
The shared attribution head L31.H21 achieves \textbf{0\%} flip rate.
Per-task patching of task-specific heads (Appendix~\ref{sec:appendix_circuits},
Figure~\ref{fig:nb18_per_task}) produces 0 verified flips for sentiment\_analysis;
translate task-specific heads reach at most 40\% (L30.H03), still below the
causal-necessity threshold: the bars are genuine partial effects, not definitive
causal bottlenecks.

\paragraph{Interpretation.}
This result is mechanistically informative.
Attribution identifies which heads \emph{contribute} to the refusal projection;
causal patching measures whether any single head is \emph{sufficient} to commit
the decision. The result is clear: refusal is \textbf{distributed}.
No single head is a bottleneck, even heads with the highest attribution or global
flip rates. Gradient-weighted attribution and causal patching thus measure different
things---prominence in the refusal direction vs.\ bottleneck status---and they
dissociate here.

This is consistent with the subspace picture: the refusal signal is spread across
many dimensions and heads; a global residual-stream intervention (Arditi) works
\emph{precisely because} it targets this distributed signal in aggregate,
while single-head patching is too local to interrupt it.
Crucially, this finding also explains why targeted head ablation (NB19, \S\ref{sec:additional})
achieves near-zero suppression: the circuit has no single-point bottleneck to remove.
Subspace-level methods (SafeConstellations) succeed by targeting the distributed
geometric structure directly.

\FloatBarrier
%% -----------------------------------------------------------
\section{Discussion}
%% -----------------------------------------------------------
\label{sec:discussion}

\paragraph{A unified account of refusal geometry.}
The results fit a coherent three-phase picture:

\begin{enumerate}
  \item \textbf{Early layers:} The harmful-refusal direction crystallises early and is
  fully task-agnostic---nearly identical across all tasks and converging before task-identity
  constellations even form. This universality is what makes global DIM ablation effective
  for harmful content.

  \item \textbf{Mid layers:} Task-identity constellations dominate representations.
  Over-refusal accumulates as a small within-cluster perturbation, geometrically distinct
  from the harmful-refusal direction, higher-dimensional, and organised around task-identity
  axes rather than a shared refusal axis.
  Crucially, a linear probe can perfectly separate over-refusal from harmful-refusal
  from the very first layer---the two phenomena are representationally distinct throughout.

  \item \textbf{Late layers:} Task-specific attention heads commit the final refusal
  decision. These heads are 80\% unique per task, confirming the task-conditioned nature
  of over-refusal. No single head is individually sufficient to flip the decision,
  consistent with distributed commitment rather than a single bottleneck.
\end{enumerate}

\paragraph{Why steering works but fails to be selective.}
Ablating the harmful-refusal direction ($\hat{\mathbf{r}}_{\text{Arditi}}$) succeeds for
harmful refusals precisely because that direction \emph{is} universal---a task-agnostic
signal that a task-agnostic intervention can cleanly target.
The over-refusal direction is not universal: it is a within-task perturbation organized
around task-identity directions.
Because the two directions share $\approx 20\%$ of variance ($\cos \approx 0.45$),
ablating $\hat{\mathbf{r}}_{\text{Arditi}}$ incidentally removes part of the over-refusal
signal, producing $-30$ pp collateral over-refusal suppression---which may look like success
if one does not also measure the cost to the safety circuit.
Measuring selectivity---over-refusal suppression per unit of safety disruption---reveals
the structural difference: 0.67 for Arditi ablation vs.\ 1.57 for SafeConstellations.

\paragraph{The selectivity framing.}
This reframes the steering problem.
The question is not ``does ablation suppress over-refusal?'' (it does, collaterally)
but ``can ablation suppress it \emph{without} corrupting the harmful-refusal circuit?''
Ablating the harmful-refusal direction cannot answer this, because that direction
straddles both phenomena---removing it disrupts both, leaving it intact fixes neither.
Task-conditioned geometry allows a different decomposition: separate the over-refusal
signal into the component shared with harmful-refusal (unavoidably disrupted) and the
component orthogonal to it (over-refusal-specific, selectively targetable).
SafeConstellations targets the second component, achieving genuine selectivity.
Its residual bypass is itself geometrically predicted: task-wrapped harmful prompts
share task-cluster identity with benign over-refusal at mid-layer, so the constellation
hook fires on them too---a consequence of targeting task-identity geometry, not a flaw.

\paragraph{Scope of the linear representation hypothesis.}
The LRH \citep{park2023linear} holds for \emph{harmful} refusal in its causal form:
one direction (the harmful-refusal direction $\hat{\mathbf{r}}_{\text{Arditi}}$) largely
suffices for ablation (80\% ASR on our evaluation set).
Importantly, this is not the same as representational one-dimensionality: harmful-refusal
variance spans 8 components at L12, with PC1 explaining only 30\%.
The causal bottleneck is one-dimensional; the representation is not.
Over-refusal breaks both senses: the over-refusal direction is more diffuse (11 components,
PC1 = 25\%) \emph{and} its principal axes are task-identity directions, so no single
direction captures it causally.

\paragraph{Calibration.}
(i) Our Arditi replication uses 25 harmful samples drawn from our evaluation set, not
JailbreakBench/AdvBench; numbers reflect in-distribution performance.
(ii) The selectivity experiment uses $n=20$ per condition; percentage-point differences
carry $\approx\pm11$ pp uncertainty (95\% CI) at this sample size.
(iii) Only 2 tasks (sentiment, translate) contribute over-refusal samples; the
0.538 pairwise cosine should be treated as preliminary evidence rather than a
definitive estimate.
(iv) 2-D PCA visualisations capture $\sim$10\% variance; all quantitative claims cite
full 4096-dim silhouette scores or high-dimensional centroid distances, not PCA-compressed
figures.
(v) Causal patching tests 5 pairs per head; rates below 50\% should be read as
indicative of distributed causal influence, not definitive causal necessity claims.

\FloatBarrier
%% -----------------------------------------------------------
\section{Additional Experiments}
\label{sec:additional}
%% -----------------------------------------------------------

\paragraph{Targeted head ablation (NB19).}
Causal patching (NB18, \S\ref{sec:causal_patching}) shows that individual head
patching cannot flip refusal decisions, suggesting distributed commitment.
We followed up with \emph{ablation} (zeroing out) of the task-specific and shared
attention heads identified in \S\ref{sec:circuits} (L27--L31) during generation,
measuring whether targeted head suppression achieves better
over-refusal/harmful-refusal selectivity than global residual-stream ablation (NB19).

\begin{table}[h]
\centering
\small
\setlength{\tabcolsep}{5pt}
\begin{tabular}{lcc}
\toprule
Condition & OR suppression & Harm suppression \\
\midrule
Shared head only (L31.H21)            & $+0$ pp & $+0$ pp \\
Task-specific heads (union, L27--L31) & $-5$ pp & $+0$ pp \\
Task-conditioned ablation             & $+0$ pp & $+5$ pp \\
\midrule
\textit{Reference: Arditi (NB15)}     & $+30$ pp & $+45$ pp \\
\textit{Reference: SafeConstel (NB15)}& $+55$ pp & $+35$ pp \\
\bottomrule
\end{tabular}
\caption{Targeted head ablation results (NB19, $n$\,=\,20 per condition). Ablating
individual heads or their union achieves negligible suppression ($\leq$5 pp), compared
to $+30$--$+55$ pp for residual-stream methods. Refusal commitment is geometrically
distributed, not head-localisable.}
\label{tab:nb19}
\end{table}

The null result confirms that refusal is \emph{distributed} in the residual stream:
no single head or task-specific head set is individually sufficient.
This rules out the hypothesis that targeted head suppression could match or surpass
global DIM ablation, and further motivates subspace-level methods (SafeConstellations)
over circuit-level interventions.

\paragraph{Cross-model replication (NB20, partial).}
We ran constellation geometry and directional analysis on pre-computed Qwen1.5-7B-Chat
embeddings (same 270 prompts, 31 layers, 4096-dim).
\textbf{Claim 1 replicates}: task-identity constellations form with inter-task
silhouette peaking at \textbf{0.458} at \textbf{L5}---comparable to LLaMA's 0.454
at L12, but concentrated earlier in the network, consistent with architectural differences
between model families.
\textbf{Claims 2--3 could not be completed}: Qwen almost never refuses on this
evaluation set (1 refused-harmful sample, 22 over-refusal vs.\ 25 and 48 for LLaMA),
leaving the harmful-refusal DIM direction undefined.
Full replication of the two-subspace finding requires an evaluation set that elicits
meaningful refusal rates from Qwen1.5-7B-Chat; we leave this to future work.

\FloatBarrier
%% -----------------------------------------------------------
\section{Conclusion}
%% -----------------------------------------------------------

Representation steering is powerful---but its effectiveness depends critically on what
is being targeted.
For harmful refusals, the target is universal: a single task-agnostic direction
mediates refusal across all tasks, which is why global DIM ablation works cleanly.
For over-refusal, the target is task-conditioned: erroneous refusals live inside
task-identity clusters, their directions vary by task, and their subspace is
higher-dimensional. These two phenomena are representationally distinct from the
first transformer layer onward.

Ablating the harmful-refusal direction to fix over-refusal is the wrong tool for
the job. The partial geometric overlap between the two directions produces collateral
suppression, which can look like success if one does not also measure the cost to
the safety circuit. Measuring selectivity reveals the structural difference: a
task-conditioned method that targets the within-task subspace directly achieves
genuine selectivity, eliminating over-refusal while preserving substantially more
of the harmful-refusal circuit.

The takeaway: the right question is not ``does steering suppress over-refusal?'' but
``can it do so \emph{selectively}?'' The answer requires a geometry-aware intervention,
not a global one.

%% -----------------------------------------------------------
\section*{Limitations}
%% -----------------------------------------------------------

Analysis is limited to LLaMA-3.1-8B-Instruct on a 270-sample evaluation set.
Over-refusal direction analysis requires sufficient over-refusal examples per task;
cryptanalysis and RAG-QA contribute zero over-refusal samples in our set, limiting
per-task analysis to sentiment ($n_{\OR}=20$) and translate ($n_{\OR}=28$).
Causal patching (NB18) uses 5 test pairs per head; flip rates should be read as
indicative of distributed causal structure rather than precise causal estimates.
Single-head patching may underestimate circuit importance when refusal arises from
distributed interactions: a head that is not individually sufficient may still
contribute jointly.
The $n=20$ per-condition selectivity experiment carries $\approx\pm11$ pp uncertainty
(95\% CI); all percentage-point differences should be interpreted accordingly.
Cross-model replication (NB20) confirms Claim 1 on Qwen1.5-7B (task-identity constellations,
silhouette 0.458 at L5) but Claims 2--3 are untestable: Qwen produces only 1
refused-harmful sample and 22 over-refusal samples on this evaluation set, leaving
directional and dimensionality findings specific to LLaMA-3.1-8B-Instruct.

%% -----------------------------------------------------------
\bibliography{custom}
%% -----------------------------------------------------------

\appendix

\section{Dataset Details}
\label{sec:appendix_data}

\textbf{Tasks:} (1) sentiment\_analysis; (2) translate; (3) cryptanalysis; (4) rag\_qa;
(5) rephrase. Task identity is determined from the \texttt{intended\_task\_labels}
field, not \texttt{text\_type\_labels}, because cryptanalysis/rag\_qa samples are
labeled \texttt{harmful\_instruction} in text-type fields even for benign phrasings.

\textbf{Sample counts.} Refused-harmful: 25; Harmless-answered: 157; Over-refusal:
48 (sentiment: 20, translate: 28, rephrase: 0, cryptanalysis: 0, rag\_qa: 0);
Target (answered benign): 169.

\section{Refusal Classification Prompt}
\label{sec:appendix_judge}

The GPT-4o refusal classification prompt used throughout (identical to original
evaluation pipeline \citep{safeconstellations2025}):

\begin{quote}
\small
\textit{There are three classes of responses:
(1)~\textbf{direct answer}---any direct attempt to give a response, regardless of accuracy;
(2)~\textbf{direct refusal}---states it will not answer, or that the question should not be answered;
(3)~\textbf{indirect refusal}---cannot answer, then offers alternatives without answering.
QUESTION: \{q\}
RESPONSE: \{r\}
Please briefly explain your reason and state your choice inside [[ and ]].
CLASS:}
\end{quote}

\section{OR Subspace Dimensionality: Layer Sweep}
\label{sec:appendix_nb16}

\begin{figure*}[h]
  \centering
  %% SOURCE: NB16 — fig_nb16_layer_sweep.png
  \includegraphics[width=0.95\linewidth]{figures/fig_nb16_layer_sweep.png}
  \caption{\textbf{Layer sweep of subspace dimensionality (NB16).}
  Left: PC1 variance fraction per layer. Right: components needed for 80\% variance.
  Harmful-refusal consistently requires fewer components from L7 onward;
  over-refusal inverts briefly at early layers before task constellations form.}
  \label{fig:nb16_layer_sweep}
\end{figure*}

\section{Extended Circuit Results}
\label{sec:appendix_circuits}

\paragraph{Causal patching summary (NB18).}
Global top heads by flip rate (5 test pairs): L13.H12, L13.H24, L13.H15, L13.H02,
L13.H08, L15.H08 --- all at 2/5 (40\%, tied maximum); L22.H19 at 1/5 (20\%);
all remaining heads including L25.H21, L31.H21, L30.H04, L30.H08 at 0/5 (0\%).

Per-task patching: \textbf{sentiment\_analysis} --- L31.H21 (shared) 0/4, L31.H30
(task-specific) 0/4, L31.H11 (task-specific) 0/4; all zero.
\textbf{translate} --- L30.H03 (task-specific) 2/5 (40\%), L31.H13 (task-specific)
1/5 (20\%), L31.H21 (shared) 1/5 (20\%).
The shared head L31.H21 achieves 0\% on sentiment but 20\% on translate,
consistent with weak cross-task rather than universal causal influence.
No head reaches the 50\% causal-necessity threshold on any task.
Full flip-rate tables for all test pairs are available on request.

\begin{figure*}[h]
  \centering
  %% SOURCE: NB18 — fig_nb18_per_task_causal.png
  %% Copy to: paper/figures/fig11_nb18_per_task_causal.png
  \includegraphics[width=0.95\linewidth]{figures/fig11_nb18_per_task_causal.png}
  \caption{\textbf{Per-task causal patching flip rates for task-specific and shared heads (NB18).}
  Blue bars: task-specific attribution heads tested on their own task.
  Purple bars ($\star$): shared head L31.H21 tested on each task.
  \textbf{sentiment\_analysis}: all heads at 0\% (0/4 pairs).
  \textbf{translate}: L30.H03 reaches 40\% (2/5 pairs), L31.H13 and L31.H21 at 20\%
  (1/5)---partial causal influence, but none reaching the 50\% necessity threshold.
  No head is individually sufficient to commit task-conditioned refusal even when
  it is the top attribution contributor for that task.
  The heads with highest global flip rates (L13 cluster at 40\%) are \emph{absent}
  from task-specific attribution sets, suggesting a weak, cross-task causal signal
  distinct from the task-conditioned attribution circuit.}
  \label{fig:nb18_per_task}
\end{figure*}

\end{document}

\end{document}
